\chapter{Probabilité conditionnelle et Bayes}\label{ch:iv1:conditional-probability-and-bayes}

\begin{ivtrains}
Mettre à jour correctement une croyance~: conditionner sur ce que l'on a réellement appris, et sur la manière dont on l'a appris. La formule de Bayes en cotes rend la plupart de ces
calculs une simple somme. Les pièges tiennent aux protocoles d'information, aux taux de base, aux observations corrélées et à la sélection, et les
questions suivantes transforment la mise à jour en un prix ou en une décision.
\end{ivtrains}

\begin{question}[$\star$ \iqroles{QR, QT, Trader, Analyste, Risque} \iqfirm{toutes} \iqlevel{junior}]\label{iq:iv1:conditional-probability-and-bayes:1}
Une famille a deux enfants, chacun ayant autant de chances d'être un garçon que d'être une fille, de façon indépendante. On vous dit qu'au moins l'un d'eux est une fille. Quelle est la probabilité que ce soient deux filles~? Et si vous rencontrez au hasard l'un des enfants et qu'il s'agit d'une fille~?
\end{question}

\begin{question}[$\star$ \iqroles{QR, QT, Trader, Analyste, DS, Risque} \iqfirm{toutes} \iqlevel{junior}]\label{iq:iv1:conditional-probability-and-bayes:2}
Une maladie touche 1~\% de la population. Un test la détecte chez 99~\% des personnes qui en sont atteintes, et il affiche à tort un résultat positif pour 5~\% des personnes qui n'en sont pas atteintes. Vous obtenez un résultat positif. Quelle est la probabilité que vous soyez malade~?
\end{question}

\begin{question}[$\star$ \iqroles{QR, QT, Trader, Analyste} \iqfirm{toutes} \iqlevel{junior}]\label{iq:iv1:conditional-probability-and-bayes:3}
Un prix se trouve derrière l'une de trois portes. Vous choisissez une porte. L'animateur, qui connaît l'emplacement du prix, ouvre toujours une porte différente, derrière laquelle il n'y a pas de prix, et vous propose de passer à la porte fermée restante. Faut-il changer de porte, et quelle est votre chance de gagner si vous le faites~? Et s'il y a 100 portes et que l'animateur en ouvre 98~?
\end{question}

\begin{question}[$\star$ \iqroles{QR, QT, Trader, Analyste, Risque} \iqfirm{toutes} \iqlevel{junior}]\label{iq:iv1:conditional-probability-and-bayes:4}
Un sac contient 100 pièces, dont 99 sont équilibrées et une a face sur ses deux faces. Vous en tirez une au hasard et obtenez 10 faces de suite en la lançant. Quelle est maintenant la probabilité que vous ayez en main la pièce à deux faces~?
\end{question}

\begin{question}[$\star$ \iqroles{QR, QT, Trader, Analyste} \iqfirm{toutes} \iqlevel{junior}]\label{iq:iv1:conditional-probability-and-bayes:5}
Trois cartes se trouvent dans un chapeau~: une rouge des deux côtés, une bleue des deux côtés, et une rouge d'un côté et bleue de l'autre. Vous en tirez une au hasard et la posez sur la table~; la face visible est rouge. Quelle est la probabilité que l'autre face soit elle aussi rouge~?
\end{question}

\begin{question}[$\star$ \iqroles{QR, QT, Analyste, DS, Risque} \iqfirm{toutes} \iqlevel{junior}]\label{iq:iv1:conditional-probability-and-bayes:6}
Dans une ville, 85~\% des taxis sont verts et 15~\% sont bleus. Après un accident, un témoin affirme que le taxi était bleu. Des tests montrent que les témoins identifient la couleur correctement 80~\% du temps, quelle que soit la couleur. Quelle est la probabilité que le taxi était bleu~?
\end{question}

\begin{question}[$\star$ \iqroles{QR, QT, Trader, Analyste} \iqfirm{toutes} \iqlevel{junior}]\label{iq:iv1:conditional-probability-and-bayes:7}
Vous lancez deux dés équilibrés, et quelqu'un les regarde et vous dit qu'au moins l'un d'eux montre un six. Quelle est la probabilité que les deux montrent un six~? Quelle est la probabilité que la somme soit égale à 7~?
\end{question}

\begin{question}[$\star$ \iqroles{QR, QT, Trader, Analyste} \iqfirm{toutes} \iqlevel{junior}]\label{iq:iv1:conditional-probability-and-bayes:8}
Un sac contient une boule, qui a autant de chances d'être rouge que bleue. Vous ajoutez une boule rouge, secouez le sac et tirez une boule, qui est rouge. Quelle est la probabilité que la boule restée dans le sac soit rouge~?
\end{question}

\begin{question}[$\star\star$ \iqroles{QR, QT, Trader} \iqfirm{hedge fund, teneur de marché} \iqlevel{junior}]\label{iq:iv1:conditional-probability-and-bayes:9}
Trois prisonniers, A, B et C, savent que l'un d'entre eux, choisi au hasard, sera gracié. A demande au gardien, qui connaît le résultat, de désigner l'un des deux autres qui ne sera pas gracié. Le gardien répond «~B~». Quelle est maintenant la probabilité que A soit gracié, et celle de C, si le gardien choisit au hasard lorsqu'il a le choix~? Et si le gardien désigne toujours B lorsque c'est possible~?
\end{question}

\begin{question}[$\star\star$ \iqroles{QR, QT, Trader} \iqfirm{hedge fund, teneur de marché} \iqlevel{junior}]\label{iq:iv1:conditional-probability-and-bayes:10}
Un prix se trouve derrière l'une de trois portes, et vous choisissez la porte 1. Première règle~: l'animateur ignore où se trouve le prix, ouvre au hasard la porte 2 ou la porte 3, et il se trouve qu'elle ne cache rien. Faut-il changer, et quel en est l'enjeu~? Seconde règle~: l'animateur le sait, n'ouvre jamais la porte du prix et, lorsqu'il a le choix, ouvre toujours la porte 2. Quelle est votre chance si vous changez après avoir vu la porte 3 ouverte~? Après avoir vu la porte 2 ouverte~?
\end{question}

\begin{question}[$\star\star$ \iqroles{QR, QT, Trader} \iqfirm{hedge fund, teneur de marché} \iqlevel{junior}]\label{iq:iv1:conditional-probability-and-bayes:11}
Une famille a deux enfants. Chacun est un garçon ou une fille avec la même chance, et chacun naît un jour quelconque de la semaine avec la même chance, de façon indépendante. Vous apprenez qu'au moins l'un des enfants est un garçon né un mardi. Quelle est la probabilité que les deux enfants soient tous deux des garçons~?
\end{question}

\begin{question}[$\star\star$ \iqroles{QR, QT, DS, Risque} \iqfirm{toutes} \iqlevel{junior}]\label{iq:iv1:conditional-probability-and-bayes:12}
Une pièce a une probabilité inconnue de faire face, et avant tout lancer vous estimez que chaque valeur entre 0 et 1 est aussi probable. Vous observez 7 faces sur 10 lancers. Quelle est la probabilité que le lancer suivant donne face~? Pourquoi pas $0.7$~?
\end{question}

\begin{question}[$\star\star$ \iqroles{QR, QT, Trader, Risque} \iqfirm{toutes} \iqlevel{junior}]\label{iq:iv1:conditional-probability-and-bayes:13}
Un sac contient deux pièces, l'une équilibrée et l'autre dont la probabilité de faire face vaut $3/4$. Vous en prenez une au hasard, la lancez deux fois et obtenez deux faces. Quelle est la probabilité que le troisième lancer de la même pièce donne lui aussi face~?
\end{question}

\begin{question}[$\star\star$ \iqroles{QR, QT, DS, Risque} \iqfirm{toutes} \iqlevel{junior}]\label{iq:iv1:conditional-probability-and-bayes:14}
Une affection a une prévalence de 1~\%. Un test donne le bon résultat dans 90~\% des cas, aussi bien chez les personnes atteintes que chez celles qui ne le sont pas, et les tests répétés sont indépendants sachant l'état réel. Vous obtenez deux résultats positifs. Quelle est la probabilité que vous soyez atteint~? Un troisième test revient ensuite négatif. Qu'en est-il maintenant~?
\end{question}

\begin{question}[$\star\star$ \iqroles{QR, QT, DS, Analyste, Risque} \iqfirm{toutes} \iqlevel{junior}]\label{iq:iv1:conditional-probability-and-bayes:15}
Sur deux ans, la stratégie A présente une proportion de transactions gagnantes supérieure à celle de la stratégie B chaque année. B peut-elle malgré tout avoir la proportion de transactions gagnantes la plus élevée sur les deux ans pris ensemble~? Donnez des chiffres si c'est possible, et dites ce que vous présenteriez.
\end{question}

\begin{question}[$\star\star$ \iqroles{QR, QT, DS, Risque} \iqfirm{hedge fund, banque} \iqlevel{junior}]\label{iq:iv1:conditional-probability-and-bayes:16}
Vous testez 1\,000 signaux de trading, dont 5~\% sont réellement efficaces. Un test au niveau de 5~\% signale à tort un signal inutile dans 5~\% des cas, et il signale un signal efficace dans 80~\% des cas. Un signal est retenu par le test. Quelle est la probabilité qu'il soit réellement efficace~? Que faudrait-il modifier en premier si cette probabilité est trop faible~?
\end{question}

\begin{question}[$\star\star$ \iqroles{QR, QT, Trader} \iqfirm{teneur de marché, hedge fund} \iqlevel{junior}]\label{iq:iv1:conditional-probability-and-bayes:17}
Une pièce a une probabilité de faire face égale à $0.6$ ou à $0.4$, et vous jugez les deux valeurs aussi probables. En 10 lancers, vous observez 7 faces et 3 piles. Quel est le prix juste d'un contrat qui verse 1~\$ si le lancer suivant donne face~?
\end{question}

\begin{question}[$\star\star$ \iqroles{QR, QT, Trader} \iqfirm{hedge fund, teneur de marché} \iqlevel{junior}]\label{iq:iv1:conditional-probability-and-bayes:18}
Une boîte contient 4 boules, et le nombre de boules rouges a autant de chances d'être 0, 1, 2, 3 ou 4. Vous en tirez deux sans remise, et les deux sont rouges. Quelle est la probabilité que la boule suivante tirée soit rouge~?
\end{question}

\begin{question}[$\star\star$ \iqroles{QR, QT, Analyste, DS, Risque} \iqfirm{toutes} \iqlevel{junior}]\label{iq:iv1:conditional-probability-and-bayes:19}
Un profil ADN correspond à une personne prise au hasard une fois sur un million. Une ville compte 5 millions d'habitants, dont l'un a commis un crime, et rien d'autre ne désigne quiconque. Une personne tirée au hasard dans la ville est testée et correspond au profil. Un avocat affirme que la probabilité d'une correspondance fortuite est d'une chance sur un million, donc que la personne est presque certainement coupable. Quelle est la probabilité de culpabilité~?
\end{question}

\begin{question}[$\star\star$ \iqroles{QT, Trader, Analyste, Risque} \iqfirm{hedge fund, banque} \iqlevel{junior}]\label{iq:iv1:conditional-probability-and-bayes:20}
Parmi l'ensemble des gérants de fonds, 10~\% sont compétents et battent le marché une année donnée avec une probabilité de 70~\%, de façon indépendante d'une année à l'autre~; les autres le font avec une probabilité de 50~\%. Un gérant a battu le marché cinq années de suite. Quelle est la probabilité que ce gérant soit compétent~? Quelle est sa chance de battre le marché l'année suivante~?
\end{question}

\begin{question}[$\star\star$ \iqroles{QT, Trader, Analyste} \iqfirm{hedge fund, teneur de marché, banque} \iqlevel{junior}]\label{iq:iv1:conditional-probability-and-bayes:21}
Une entreprise vaut un montant $V$ pour son propriétaire, uniformément réparti entre 0 et 100. Seul le propriétaire connaît $V$, et l'entreprise vaudrait $1.5\,V$ sous votre gestion. Vous formulez une seule offre $b$, et le propriétaire vend si $b \ge V$. Quelle offre maximise l'espérance du profit, et quelle est cette espérance~? Et si l'entreprise valait $2.5\,V$ pour vous~?
\end{question}

\begin{question}[$\star\star$ \iqroles{QR, QT, Trader} \iqfirm{hedge fund, teneur de marché} \iqlevel{junior}]\label{iq:iv1:conditional-probability-and-bayes:22}
Une urne contient au départ une boule rouge et une boule bleue. À chaque tirage, vous prenez une boule au hasard, puis vous la remettez avec une boule supplémentaire de même couleur. Quelle est la probabilité que les 10 premiers tirages soient tous rouges~? Sachant que le dixième tirage est rouge, quelle est la probabilité que le premier tirage ait donné une boule rouge~?
\end{question}

\begin{question}[$\star\star$ \iqroles{QR, QT, DS, Risque} \iqfirm{toutes} \iqlevel{junior}]\label{iq:iv1:conditional-probability-and-bayes:23}
Une pièce a une probabilité de faire face $p$ tirée uniformément entre 0 et 1. Elle est ensuite lancée deux fois. Les deux lancers sont-ils indépendants~? Quelle est la probabilité d'obtenir deux faces, et quelle est la corrélation entre les deux résultats~?
\end{question}

\begin{question}[$\star\star$ \iqroles{QR, QT, Risque, DS} \iqfirm{toutes} \iqlevel{junior}]\label{iq:iv1:conditional-probability-and-bayes:24}
$X$ et $Y$ sont conjointement gaussiens, chacun de moyenne 0 et de variance 1, et de corrélation $0.5$. Quelles sont la moyenne et la variance de $Y$ sachant que $X = 2$~? Quelle est la probabilité que $Y > 0$ sachant que $X > 0$~?
\end{question}

\begin{question}[$\star\star$ \iqroles{QR, QT, Risque, DS} \iqfirm{hedge fund, banque} \iqlevel{junior}]\label{iq:iv1:conditional-probability-and-bayes:25}
Votre croyance sur la fair value d'une action est une loi normale, centrée sur 100, d'écart-type 4. Un signal bruité affiche ensuite 106, le signal étant la fair value augmentée d'un bruit normal indépendant d'écart-type 3. Quelle est votre moyenne mise à jour, et quel est votre écart-type~?
\end{question}

\begin{question}[$\star\star$ \iqroles{QT, Trader, QR} \iqfirm{teneur de marché, crypto} \iqlevel{junior}]\label{iq:iv1:conditional-probability-and-bayes:26}
Un titre vaudra 90 ou 110, avec la même probabilité. Un trader arrive et échange avec vous une action. Avec une probabilité de 20~\%, le trader est un trader informé~: il achète lorsque la vraie valeur vaut 110 et vend lorsqu'elle vaut 90~; sinon, il achète ou vend avec la même probabilité. Quel est le prix le plus bas auquel vous pouvez proposer de vendre une action sans perdre en moyenne~? Après une telle vente, où ce prix doit-il se placer pour le trader suivant~?
\end{question}

\begin{question}[$\star\star$ \iqroles{QR, QT, Trader, Analyste} \iqfirm{hedge fund, banque} \iqlevel{junior}]\label{iq:iv1:conditional-probability-and-bayes:27}
Deux analystes prévoient chacun correctement la direction du marché dans 70~\% des cas, et la hausse et la baisse sont équiprobables. Tous deux annoncent la hausse. Quelle est la probabilité que le marché monte si leurs erreurs sont indépendantes~? Et s'ils sont toujours d'accord, parce qu'ils lisent le même rapport~? Et s'ils ont tous deux raison dans 55~\% des journées et tous deux tort dans 15~\%~?
\end{question}

\begin{question}[$\star\star\star$ \iqroles{QR, QT, Trader} \iqfirm{hedge fund, teneur de marché} \iqlevel{junior, confirmé}]\label{iq:iv1:conditional-probability-and-bayes:28}
Vous recevez 13 cartes d'un jeu de 52 cartes. Sachant que vous avez au moins un as, quelle est la probabilité d'en avoir au moins deux~? Sachant que vous avez l'as de pique, quelle est cette probabilité~? Pourquoi ces deux probabilités sont-elles différentes~?
\end{question}

\begin{question}[$\star\star\star$ \iqroles{QR, QT, Trader} \iqfirm{hedge fund, teneur de marché} \iqlevel{junior, confirmé}]\label{iq:iv1:conditional-probability-and-bayes:29}
Un prix se cache derrière l'une de quatre portes, et vous en choisissez une. Deux fois de suite, l'animateur, qui sait où se trouve le prix, ouvre une porte qui n'est ni celle que vous avez choisie ni celle qui cache le prix, au hasard lorsqu'il y a un choix. Après chaque ouverture, vous pouvez garder votre porte ou passer à une autre porte fermée, au hasard s'il y en a plusieurs. Après la seconde ouverture, la porte que vous détenez alors est ouverte. Quelle est la meilleure suite de décisions, et quelle est votre chance de gagner~?
\end{question}

\begin{question}[$\star\star\star$ \iqroles{QR, QT, DS, Risque} \iqfirm{hedge fund, banque} \iqlevel{junior, confirmé}]\label{iq:iv1:conditional-probability-and-bayes:30}
Dans un pays, chaque famille a des enfants jusqu'au premier garçon, puis s'arrête. Chaque enfant est un garçon ou une fille avec la même probabilité. Quel est le nombre moyen d'enfants par famille~? Quelle est la proportion de garçons parmi les enfants~? Vous rencontrez au hasard un enfant de toute la population~: en moyenne, combien d'enfants, celui-ci compris, compte la famille de cet enfant~?
\end{question}

\begin{question}[$\star\star\star$ \iqroles{QR, QT, DS, Analyste} \iqfirm{hedge fund, banque} \iqlevel{junior, confirmé}]\label{iq:iv1:conditional-probability-and-bayes:31}
Chaque candidat a deux notes indépendantes, $X$ et $Y$, uniformes entre 0 et 1, et n'est embauché que si $X + Y > 1$. Quelle est la corrélation entre $X$ et $Y$ parmi les personnes embauchées~? Comment cela pourrait-il induire en erreur quelqu'un qui n'étudie que ces personnes~?
\end{question}

\begin{question}[$\star\star\star$ \iqroles{QR, QT, Risque} \iqfirm{hedge fund, banque} \iqlevel{junior, confirmé}]\label{iq:iv1:conditional-probability-and-bayes:32}
Un marché se trouve dans l'un de deux régimes~: calme, avec des rendements quotidiens normaux de moyenne 0 et de volatilité 1~\%, ou de crise, avec une moyenne nulle et une volatilité de 3~\%. Votre croyance aujourd'hui est de 5~\% de crise. Le rendement d'aujourd'hui est de $-4~\%$. Quelle est votre probabilité mise à jour que le marché soit en crise~? Et si le rendement d'aujourd'hui est exactement 0~?
\end{question}

\begin{question}[$\star\star\star$ \iqroles{QR, QT, DS} \iqfirm{hedge fund, teneur de marché} \iqlevel{junior, confirmé}]\label{iq:iv1:conditional-probability-and-bayes:33}
$X$ et $Y$ sont des variables aléatoires exponentielles indépendantes de moyenne 1. Sachant que $X + Y = s$, quelle est la loi de $X$~? Comment cela change-t-il si $X$ et $Y$ sont plutôt des normales centrées réduites indépendantes~?
\end{question}

\begin{question}[$\star\star\star$ \iqroles{QR, QT, Trader} \iqfirm{hedge fund, teneur de marché} \iqlevel{junior, confirmé}]\label{iq:iv1:conditional-probability-and-bayes:34}
Deux enveloppes contiennent de l'argent, l'une ayant deux fois plus que l'autre, et vous en prenez une au hasard. Il semble qu'avec un montant $x$ dans votre enveloppe, l'autre contienne $2x$ ou $x/2$ avec la même probabilité, soit en moyenne $1.25x$, donc vous devriez toujours changer, puis changer de nouveau. Où est l'erreur~? Montrez-le à l'aide de montants tirés des couples $(1, 2)$, $(2, 4)$ et $(4, 8)$, chaque couple étant équiprobable.
\end{question}

\begin{question}[$\star\star\star$ \iqroles{QR, QT, Risque, DS} \iqfirm{toutes} \iqlevel{junior, confirmé}]\label{iq:iv1:conditional-probability-and-bayes:35}
Un bus circule selon un horaire fixe, soit toutes les 15 minutes, soit toutes les 30 minutes, et vous estimez que les deux sont également probables. Vous arrivez à un instant quelconque et attendez 10 minutes sans voir de bus. Quelle est la probabilité qu'il circule toutes les 30 minutes~? En moyenne, combien de temps encore attendrez-vous~?
\end{question}

\begin{question}[$\star\star\star$ \iqroles{QR, QT, Trader} \iqfirm{hedge fund, teneur de marché} \iqlevel{junior, confirmé}]\label{iq:iv1:conditional-probability-and-bayes:36}
Une urne est de type A ou de type B, avec la même probabilité. Des personnes devinent le type l'une après l'autre, et chaque supposition est annoncée. Avant de deviner, chaque personne observe un signal privé, qui désigne le vrai type avec une probabilité de $2/3$, indépendamment des autres signaux. Chaque personne entend aussi toutes les suppositions précédentes, mais ne voit jamais les signaux dont elles découlent. Chacune devine le type le plus probable compte tenu de tout cela, et suit son propre signal si les deux types sont équiprobables. Finalement, la supposition d'une personne ne dépend plus de son propre signal, et à partir de là tout le monde devine le même type. Quelle est la probabilité que ce type soit le mauvais~? En moyenne, combien de personnes devinent avant la première dont le signal n'a plus d'importance~?
\end{question}

\begin{question}[$\star\star\star$ \iqroles{QR, QT, Trader} \iqfirm{teneur de marché, hedge fund} \iqlevel{junior, confirmé}]\label{iq:iv1:conditional-probability-and-bayes:37}
Une pièce a une probabilité de $0.6$ ou de $0.4$ de tomber sur face, et vous estimez que les deux possibilités sont équiprobables. Vous pariez 1~\$ à parité sur trois lancers successifs, en choisissant face ou pile avant chaque lancer et en voyant tous les résultats précédents. Quelle est votre meilleure espérance de gain~? Quelle serait-elle si vous connaissiez le type de la pièce dès le départ~?
\end{question}

\begin{question}[$\star\star\star\star$ \iqroles{QR, QT, DS} \iqfirm{hedge fund, teneur de marché} \iqlevel{confirmé}]\label{iq:iv1:conditional-probability-and-bayes:38}
Vous lancez une pièce 100 fois et obtenez 60 faces. Un test bilatéral donne une p-valeur voisine de 5~\%. Vous pensez aussi qu'avant d'avoir vu les données, la pièce avait 50~\% de chances d'être parfaitement équilibrée, et que sinon sa probabilité de face était uniforme entre 0 et 1. Quelle est votre probabilité mise à jour que la pièce soit équilibrée~? Et pour 5\,100 faces sur 10\,000 lancers, qui donne une p-valeur voisine~? Comment peut-on concilier ces deux résultats~?
\end{question}

\begin{question}[$\star\star\star\star$ \iqroles{QR, QT, DS, Risque} \iqfirm{hedge fund} \iqlevel{confirmé}]\label{iq:iv1:conditional-probability-and-bayes:39}
Cent gérants ont chacun un ratio de Sharpe annuel vrai tiré d'une loi normale centrée en 0, d'écart-type $0.5$. Vous mesurez chacun avec un bruit normal indépendant d'écart-type $0.5$. En moyenne, quel ratio de Sharpe mesuré le meilleur d'entre eux affiche-t-il~? Quelle est son espérance de ratio de Sharpe vrai~? Que deviennent ces deux chiffres si aucun gérant n'a la moindre compétence~?
\end{question}

\begin{question}[$\star\star\star\star$ \iqroles{QR, QT, Trader} \iqfirm{hedge fund, teneur de marché} \iqlevel{confirmé}]\label{iq:iv1:conditional-probability-and-bayes:40}
Quelqu'un écrit deux nombres différents sur deux cartes, par la méthode de son choix, et vous ignorez tout de la méthode employée. Vous retournez une carte, choisie au hasard, et voyez son nombre. Vous dites ensuite si l'autre nombre est plus grand ou plus petit. Pouvez-vous garantir une probabilité d'avoir raison strictement supérieure à un demi, quels que soient les deux nombres~? De combien cette probabilité est-elle supérieure à un demi si ces nombres sont 10 et 20 et que vous utilisez un seuil aléatoire normal de moyenne 15 et d'écart-type 10~?
\end{question}
